\section{Conclusion}
\label{sec:conclusion}

Large language models are becoming working infrastructure for air pollution
policy, a domain where an inaccurate answer about a binding standard, a measured
concentration, or an operational response carries direct public-health stakes
in the Global South, where the PM\textsubscript{2.5} mortality burden is
heaviest~\citep{gbd2019risk,requia2024shortterm}. We measured geographic and
persona-modulated bias in that domain with a 25-country, 14-model,
four-language benchmark (English plus the native language of nine countries),
scored against ground truth anchored to official primary sources, by code
wherever the answer is a register value and by the mean of a three-vendor judge
panel elsewhere.

None of the remedies an agency can try works, and the most intuitive one
backfires. Asking in the country's own language lowers accuracy by
$\nnativa$~percentage points, most sharply for Hindi, the language least
represented in training corpora, and the penalty is not an artefact of unanswered
prompts. A local-manager persona is equivalent to no persona within two
percentage points. Regional instruction tuning at 7B returns the value in force
no more often than the base model it was tuned from. Accuracy is lowest on the
recall of a binding value, the task a manager can least afford to have wrong,
and a third of the wrong answers cite a rung of the country's own official
ladder rather than an invented number. A Global North/South gap persists on the
outcomes no judge touches, where Global South countries return the register
value $\ncodegap$~pp less often with the country as the unit. Its concentration
on tasks whose answer is a national fact is a pattern of point estimates:
every per-task interval at the country level includes~1. The development
gradient does not reach significance and its interval includes both zero and the
pre-specified threshold. On mechanism, what predicts the country-specificity
deficit at the country level is development rather than encyclopedic coverage,
and the two are too collinear in 25 countries to be told apart; the
language-corpus channel is excluded.

For Global South environmental agencies the practical message is cautionary.
LLM outputs on binding standards and measured data should be treated as drafts
checked against the official gazette, and none of the cheap prompt- or
model-level workarounds substitutes for that verification, least of all
switching to the local language, which on this evidence makes the answer worse.
The penalty falls on the languages under-served in training
corpora~\citep{joshi2020state}, hardest on Hindi; whether those languages are
also spoken where the health burden of air pollution is greatest is not
something this design measures, and the one country-level observation we have
(India) points the other way (Section~\ref{sec:discussion}).

One methodological result generalises beyond the domain. Rebuilding two task
keys from official registers and moving range comparison out of the judge and
into code changed the substantive conclusions on the same responses, and
carrying inference at the country rather than the response level changed the
status of the mechanism test. When an effect in this literature proves
fragile, the ground-truth key and the unit of inference deserve the audit
before the sample does, because a defective key adds structured noise that
concentrates where official data are hardest to obtain, and a response-level
$p$-value overstates what a handful of countries can show.

Protocol, prompts, official-source ground-truth registries, analysis code, and
anonymised response data are released openly, with a one-command reproduction
pipeline and a method-audit gate that recomputes every reported number from the
raw data. The study also provides a reusable template for applied LLM-equity
audits, with official-source ground truth, code adjudication wherever the
answer is a checkable number, and end-to-end reproducibility, in the
high-stakes public-sector domains where these systems are now being deployed.
